\section{Final Predictive Relations per Method}
\label{sec:laws}

\begin{figure}[t]
\centering
\includegraphics[width=0.56\textwidth]{figs/predictive_responses.pdf}
\caption{Final predictive relations and measured responses on each arm's native configuration axis. Left: pruning of Pythia-1.4B at step 112k (measured points) with the frozen power form of Eq.~\ref{eq:power} and the frozen median development curve $\widetilde f_c(d)$; the shaded range is covered by the 17-state refit. Middle: grouped quantization on two development states, measured math response against group size at three bit-widths, with the frozen same-input piecewise interpolation (solid, extended to the never-measured $g=32$ and $512$) and the per-configuration median (dashed). Right: the twelve $U{=}200$ distillation trajectories against reuse count with the frozen selected forms and the frozen strongest baseline. All lines were frozen before the plotted points were measured.}
\label{fig:responses}
\end{figure}

Each arm delivers one relation in the interface of Eq.~\ref{eq:interface};
we state the relation, its inputs and range, how it was selected, and
where it stops. Two roles are kept apart throughout: the reference
anchor $L_{0,c}$ is a measured quantity needed to reconstruct an
absolute loss, and the response relation $\widehat\Delta_{m,c}$ is the
predictive claim; a good absolute prediction with a source-free
$\widehat\Delta$ says nothing about whether $N_0$ or $D_0$ predict the
damage.

\subsection{Conditioning on the Initial Model State}
\label{sec:conditioning}

On the nine-state development panel with one stage held out, a
per-setting regression of the pruning response on $(N_0,D_0,L_{0,c})$
lowers the error from 0.413, 0.516, and 0.500 nats to 0.253, 0.292,
and 0.806 for math, code, and QA (interval above zero for code,
crossing for math, negative for QA); adding $L_{0,c}$ to $(N_0,D_0)$
has intervals above zero for math ($+0.55$, $[0.29,1.00]$) and code
($+0.66$, $[0.38,1.13]$), adding $D_0$ does not (Fig.~\ref{fig:inputs}A).
The source-conditioned forms below therefore use the three inputs per
capability, with the reference anchor as the carrier, and every arm
also carries a source-free relation for the cases where the source
inputs add nothing.

\subsection{Pruning}
\label{sec:law_prune}

\paragraph{Relation.} For a seen size at an unseen density in the
development range, the compact form
\begin{equation}
\widehat{\Delta L}_c(\mathbf x,d)=A_c(\mathbf x)\,\Big(\frac{1-d}{0.3}\Big)^{\gamma_c},
\qquad
A_c(\mathbf x)=\beta_c\cdot\big(1,\tilde{\log N_0},\tilde L_{0,c},\tilde{\log D_0}\big),
\label{eq:power}
\end{equation}
with five parameters per capability (tildes standardized on the
development set) approximates the twenty-parameter per-density
regression A2; for a new size or stage, and for QA everywhere, the
delivered relation is the source-free median development curve,
$\widehat L_c=L_{0,c}+\widetilde f_c(d)$.
\paragraph{Inputs and range.} $N_0$, $D_0$, and $L_{0,c}$ for
Eq.~\ref{eq:power}; $L_{0,c}$ alone for the median curve. The
nine-state freeze covers $d\in[0.6,0.9]$ and $0.55$ was an
extrapolation for it; the 17-state refit used for the confirmations
includes 0.55.
\paragraph{Selection.} Fit on nine states at $d\in\{0.9,0.8,0.7,0.6\}$
and scored at the unseen 0.65 and 0.55, Eq.~\ref{eq:power} cuts the
strength-only error two- to threefold and A2 matches it
(Table~\ref{tab:main_final}); refit on 17 states under a pre-committed
leave-one-source-out rule, A2 (0.479) and the power form (0.494) tie
within 0.02 nats and the rule prefers the continuous form that
vanishes at $d=1$. Its coefficients are in Appendix~\ref{sec:app_coef};
$\gamma_c$ is a fitted constant whose value these data do not
establish as necessary.
\paragraph{Tests and boundary.} On three checkpoints that entered no
fit the power form reaches, with fewer parameters, an error close to
A2 on math (0.243 against 0.230) and does not beat the median curve on
code (0.236 against 0.214) or QA (0.68 against 0.22); on a 2.8B source
the median curve is three to seven times better than either source
form (\S\ref{sec:unseen_settings}, \S\ref{sec:unseen_states}). The
compact form is therefore delivered as an approximation with a
tested range, and the median curve as the relation for new states.

\subsection{Quantization}
\label{sec:law_quant}

\paragraph{Relation.} Per-channel RTN offers no continuous axis: int8
and int6 stay within 0.01 nats of the reference, int4 costs 0.03 to
0.35 nats, and int3 collapses most models; at a seen bit-width a new
state is predicted by the per-bit source regression of
\S\ref{sec:conditioning} (four coefficients per bit), which improves on
the per-bit median for math and code and not for QA. Sharing one scale
across $g$ contiguous weights adds a second axis. The delivered grouped
relation for math and code on seen states is a same-input piecewise
interpolation: a per-configuration regression on
$(N_0,D_0,L_{0,c})$ at each measured $(b,g)$ cell, interpolated
linearly in $\log_2 g$ between measured group sizes at the same
bit-width and extrapolated linearly beyond the grid with a zero floor;
for QA, and for every capability on a new state, the delivered
relation is the per-configuration median of the development responses
under the same boundary rule.
\paragraph{Inputs and range.} The measured development grid
($b\in\{3,4,5\}$, $g\in\{64,128,256\}$ on six states) plus
$(N_0,D_0,L_{0,c})$ for the interpolation; $L_{0,c}$ alone for the
median.
\paragraph{Selection.} Two candidates with parametric structure were
tested first: a separable amplitude-times-shape form and a low-order
surface in $\log_2 q_{\max}$ and $\log_2 g$. The surface beat zero
change at the unseen $b=4$ and $g=128$ on the two-bit-level development
set, but its quadratic term is unidentified there (rank 16 of 20; the
curvature is set by the ridge penalty, Appendix~\ref{sec:app_forms}),
and in a confirmation round with all 54 unblinded cells as development
set (full rank) and 21 cells frozen at unseen group sizes and on a new
state, the development leave-one-state-out rule selected the surface
for math, the median for code, and zero for QA; the interpolation won
the frozen cells for math and code and the median won the new state,
so the delivered rule was fixed after that round and is labeled
retrospective (Appendix~\ref{sec:app_results}).
\paragraph{Tests and boundary.} At $g\in\{32,512\}$ on the development
states the interpolation has MAE 0.07, 0.12, and 0.48 against 0.22,
0.31, and 1.23 for the surface; on the new state the per-configuration
median is best for math and QA at the new group sizes (0.13 and 0.13)
while the interpolation is better for code (0.17 against 0.22) and zero
change is better for QA at seen configurations (0.05 against 0.16), so
source conditioning adds nothing there and the parametric surface is
not delivered (\S\ref{sec:unseen_settings}).

\subsection{Distillation}
\label{sec:law_distill}

\paragraph{Relation.} The response is the signed transfer $\delta_c$ of
an existing student relative to its own reference loss; the teacher
never enters it, and the coordinates are the supervised completion
tokens $T$, the pool size $D_U$, and the reuse count $E=T/D_U$. The
delivered relations, frozen for the confirmation, are
\begin{equation}
\widehat\delta_c=a_c\log(1+E)\ \ (\text{math, code}),\qquad
\widehat\delta_{\mathrm{QA}}=u\,(a+b\,u+q\,v),\ \ u=\log(1+T/T_{\mathrm{ref}}),\ v=\log(D_U/D_{\mathrm{ref}}),
\label{eq:kdforms}
\end{equation}
one parameter for math and code, which vanish at zero budget, and
three for QA. QA improves under every recipe we ran while math and
code move by a few hundredths to a few tenths of a nat, and with a
small pool the QA loss rises again at large $E$.
\paragraph{Inputs and range.} $T$ and $D_U$, and the student's
reference loss; no student descriptor. The relations hold on the two
development students (270M, 1B); QA on the 2Wiki distribution only,
where the gain replicates on a fresh sample and reverses sign on
MuSiQue and TriviaQA (Appendix~\ref{sec:app_measure}).
\paragraph{Selection.} Eight forms in $T$, $E$, and a source term were
frozen on twelve development runs and scored on nine unseen-pool
trajectories with the headline chosen after the test; two
descriptor-carrying forms were compared retrospectively
(Appendix~\ref{sec:app_forms}). The delivered forms were then
re-selected on all 25 trajectories by a pre-committed
leave-one-trajectory-out rule (ties within 0.02 nats toward fewer
parameters) and frozen for six new pools of $U{=}200$ at three new
budgets. On the nine Pythia states as students, a linear form in the
student state halves the constant's error for math (0.014 against
0.031) and is worse for QA; it serves only as the math predictor for
distillation candidates in \S\ref{sec:selection}.
\paragraph{Tests and boundary.} At the new pool size the selected form
beats the frozen strongest baseline for code on both students
(0.004 and 0.005 nats, intervals above zero), where the baseline is
the same $\log(1+E)$ form with an intercept, so the gain supports the
zero-budget boundary and parsimony of the form and no full
student-data-training law; for QA the joint form gains on 270M and
not on 1B, and for math a budget-only or surface baseline is better on
both students. A student 6.6 times the reference size is outside the
range: the frozen constant is better there at an unseen pool
(\S\ref{sec:unseen_states}).
